\editionPage[\pdfbookmark[0]{Fiches outils}{bm-outils-debut-parent}]{outils-debut}{Fiche outil}{Calculer et contrôler}{À consulter au fil des chapitres.}
\begin{blocEdition}{outils-debut}
\editionEtape{Poser}Aligner les virgules pour additionner ou soustraire ; compléter par des zéros. Pour multiplier, calculer sans les virgules, puis placer autant de décimales que dans les deux facteurs réunis.\editionEtape{Diviser}Division euclidienne : $a=b\times q+r$ avec $r<b$. Division décimale : transformer chaque reste en dixièmes, puis en centièmes ; arrêter quand le reste est nul ou donner un arrondi.\editionEtape{Par 10, 100, 1 000}Multiplier fait avancer chaque chiffre de 1, 2 ou 3 rangs ; diviser, ou multiplier par $0{,}1$, $0{,}01$, $0{,}001$, le fait reculer.\editionEtape{Contrôler}Estimer un ordre de grandeur avant, vérifier après : une soustraction par une addition, une division par une multiplication.\editionRetenir{Un calcul n’est terminé qu’après son contrôle : un résultat loin de l’estimation se refait.}
\end{blocEdition}
\editionPage{outils-2}{Fiche outil}{Fractions et pourcentages}{À consulter au fil des chapitres.}
\begin{blocEdition}{outils-2}
\editionEtape{Égalités}On ne change pas une fraction en multipliant, ou en divisant, numérateur et dénominateur par un même nombre non nul.\editionEtape{Comparer}Même dénominateur : comparer les numérateurs. Sinon, écrire les deux fractions avec un même dénominateur.\editionEtape{Calculer}Additionner ou soustraire : même dénominateur, puis opérer sur les numérateurs. Multiplier par un entier : multiplier le numérateur. Prendre $\frac ab$ d’un nombre : diviser par $b$, multiplier par $a$.\editionEtape{Pourcentages}$p$~\% $=\frac p{100}$. 50~\% est la moitié, 25~\% le quart, 10~\% le dixième. Pour appliquer un pourcentage, multiplier par $\frac p{100}$.\editionRetenir{À savoir par cœur : $\frac14=0{,}25$ ; $\frac12=0{,}5$ ; $\frac34=0{,}75$ ; $\frac32=1{,}5$ ; $\frac52=2{,}5$.}
\end{blocEdition}
\editionPage{outils-3}{Fiche outil}{Construire avec les instruments}{À consulter au fil des chapitres.}
\begin{blocEdition}{outils-3}
\editionEtape{Notations}Droite (AB), segment [AB], demi-droite [AB), longueur AB, angle $\widehat{ABC}$ de sommet B.\editionEtape{Instruments}Règle pour tracer ; règle graduée pour mesurer ; équerre pour les angles droits ; compas pour les cercles et le report de longueurs ; rapporteur pour mesurer et construire les angles.\editionEtape{Constructions clés}Milieu ; médiatrice au compas ; bissectrice au compas ; triangle de trois longueurs avec deux arcs ; symétrique d’un point avec l’équerre et le compas ; cercle circonscrit avec deux médiatrices.\editionEtape{Coder}Un même signe pour des longueurs égales, un petit carré pour un angle droit, des arcs identiques pour des angles égaux.\editionRetenir{Une figure à main levée, codée, aide à préparer une construction précise.}
\end{blocEdition}
\editionPage{outils-4}{Fiche outil}{Unités et conversions}{À consulter au fil des chapitres.}
\begin{blocEdition}{outils-4}
\editionEtape{Longueurs}1 km $=1\,000$ m ; 1 hm $=100$ m ; 1 dam $=10$ m ; 1 m $=10$ dm $=100$ cm $=1\,000$ mm.\editionEtape{Aires}$1\ \mathrm m^2=100\ \mathrm{dm}^2$ ; $1\ \mathrm{dm}^2=100\ \mathrm{cm}^2$ ; donc $1\ \mathrm{cm}^2=0{,}01\ \mathrm{dm}^2$ et $1\ \mathrm{dm}^2=0{,}01\ \mathrm m^2$.\editionEtape{Volumes}1 cm$^3$ est le volume d’un cube de 1 cm d’arête.\editionEtape{Durées}1 j $=24$ h ; 1 h $=60$ min $=3\,600$ s ; 1 min $=60$ s. $0{,}5$ h $=30$ min ; $0{,}25$ h $=15$ min ; $0{,}75$ h $=45$ min ; $0{,}1$ h $=6$ min. Une année : 365 jours, 366 si elle est bissextile.\editionRetenir{Pour les longueurs, deux unités voisines sont dans un rapport de 10 ; pour les aires, de 100.}
\end{blocEdition}
